% !TEX root = task-1-9-2.tex
% Editable task bank. Append new BankTask blocks at the end of the document.
\documentclass[12pt,a4paper]{article}
\usepackage{physics-textbook}
\usepackage{graphicx,needspace}
\usetikzlibrary{decorations.pathmorphing}
\newif\ifShowSolutions
\ifdefined\TutorEdition
  \ShowSolutionstrue
\else
  \ShowSolutionsfalse % Change to \ShowSolutionstrue to display solutions.
\fi
\newcounter{banktask}
\newsavebox{\BankTaskBox}
\newcommand{\BankTaskContent}[3]{%
  \noindent\textbf{\thebanktask.}\quad #2\par
  \ifShowSolutions
    \Needspace{50mm}% Keep the solution heading with its opening text or diagram.
    #3\par
    \ifdefined\TutorEdition
      \ifcsname TutorSource#1\endcsname
        \nopagebreak\smallskip{\small\csname TutorSource#1\endcsname\par}%
      \fi
    \fi
  \fi
}
\newcommand{\BankTask}[3]{%
  \refstepcounter{banktask}\par\medskip
  \sbox{\BankTaskBox}{\begin{minipage}{\linewidth}%
    \BankTaskContent{#1}{#2}{#3}\end{minipage}}%
  \ifdim\dimexpr\ht\BankTaskBox+\dp\BankTaskBox\relax>\dimexpr\textheight-20mm\relax
    \Needspace{35mm}\BankTaskContent{#1}{#2}{#3}%
  \else
    \Needspace{\dimexpr\ht\BankTaskBox+\dp\BankTaskBox\relax}%
    \noindent\usebox{\BankTaskBox}\par
  \fi
  \medskip
}
\begin{document}
\BookSubsection{1.9.2}{Закон Гука и упругая деформация}


% Source sdamgia-302: https://phys-ege.sdamgia.ru/problem?id=302
\BankTask{sdamgia-302}{
На сколько сантиметров растянется пружина, жесткость которой $k=1000\,\text{Н}/\text{м},$ под действием силы $100\,\text{Н}$? Пружину считайте идеальной.
}{
\smallskip\noindent\textbf{Решение}\par
Удлинение пружины жесткости $k=1000\,\text{Н}/\text{м}$ под действием силы $F=100\,\text{Н}$ можно определить по формуле:

\[\Delta x= \frac{F}{k} = \frac{100\,\text{Н}}{1000\,\text{Н}/\text{м}} =0{,}1\,\text{м}=10\,\text{см}.\]

Ответ: $10$.
}

% Source sdamgia-328: https://phys-ege.sdamgia.ru/problem?id=328
\BankTask{sdamgia-328}{
Под действием силы $4{,}5\,\text{Н}$ пружина удлинилась на $6\,\text{см}$. Чему равен модуль силы, под действием которой удлинение этой пружины составит $4\,\text{см}$? (В ньютонах.)
}{
\smallskip\noindent\textbf{Решение}\par
Согласно закону Гука, удлинение пружины пропорционально растягивающей ее силе: $F_{\text{упр}}=k\Delta x.$ Таким образом, потребуется сила:

\[\frac{4\,\text{см}}{6\,\text{см}} \cdot  4{,}5\,\text{Н}=3\,\text{Н},\]

чтобы растянуть пружину на $4\,\text{см}$.

Ответ: $3$.
}

% Source sdamgia-333: https://phys-ege.sdamgia.ru/problem?id=333
\BankTask{sdamgia-333}{
Две пружины растягиваются одинаковыми силами $F$. Жесткость первой пружины $k_1$ в $1{,}5$ раза больше жесткости второй пружины $k_2.$ Чему равно отношение удлинений пружин $\frac{\Delta l_2}{\Delta l_1}$?
}{
\smallskip\noindent\textbf{Решение}\par
Согласно закону Гука, растяжение пружины связано с коэффициентом жесткости и растягивающей силой соотношением $F=k\Delta l.$ Поскольку пружины растягивают одинаковыми силами, получаем:

\[k_1\Delta l_1=k_2\Delta l_2 \Longleftrightarrow  \frac{\Delta l_2}{\Delta l_1} = \frac{k_1}{k_2} =1{,}5.\]

Ответ: $1{,}5$.
}

% Source sdamgia-32578: https://phys-ege.sdamgia.ru/problem?id=32578
\BankTask{sdamgia-32578}{
На легкой пружине подвешен груз массой $2\,\text{кг}$. В состоянии покоя он растягивает пружину на $4\,\text{см}$. Какой груз следует дополнительно подвесить к этой пружине, чтобы ее удлинение увеличилось на $5\,\text{см}$? (В килограммах.)
}{
\smallskip\noindent\textbf{Решение}\par
При подвешивании груза, он находится в равновесии, поэтому сила упругости пружины уравновешивается силой тяжести груза. По закону Гука сила упругости равна $F_{\text{упр}}=kx.$

Тогда получаем соотношение: $mg=kx.$ Из данного соотношения составим пропорцию:

\[\frac{m_1}{m_1 + m_2} = \frac{x_1}{x_1 + x_2} ,\]

откуда находим массу дополнительного груза:

\[m_2= \frac{m_1x_2}{x_1} = \frac{2 \cdot  5}{4} =2{,}5\,\text{кг}.\]

Ответ: $2{,}5$.
}

% Source sdamgia-38042: https://phys-ege.sdamgia.ru/problem?id=38042
\BankTask{sdamgia-38042}{
Две лёгкие пружины прикреплены к потолку — каждая одним из своих концов. Для растяжения первой пружины на небольшую длину $\Delta l$ требуется приложить силу, равную по модулю $10\,\text{Н}$. Для растяжения второй пружины на такую же длину $\Delta l$ требуется сила, равная по модулю $4\,\text{Н}$. В обоих случаях силы прикладываются к свободным концам пружин. Определите отношение жёсткостей $\frac{k_1}{k_2}$ первой и второй пружин.
}{
\smallskip\noindent\textbf{Решение}\par
Из закона Гука $F_{\text{упр}}=k\Delta l.$ Тогда отношение жесткостей пружин при растяжении на одну и ту же длину равно

\[\frac{k_1}{k_2} = \frac{F_1}{F_2} = \frac{10}{4} =2{,}5.\]

Ответ: $2{,}5$.
}

% Source sdamgia-38487: https://phys-ege.sdamgia.ru/problem?id=38487
\BankTask{sdamgia-38487}{
В таблице приведена зависимость силы упругости от удлинения пружины. Определите жесткость пружины. (В ньютонах на метр.)

 




\begin{center}\begin{tabular}{|c|c|c|c|c|c|}\hline
$F_{\text{упр}}$, $\text{Н}$ & $4$ & $6$ & $8$ & $10$ & $12$\\\hline
$x$, $\text{м}$ & $0{,}04$ & $0{,}06$ & $0{,}08$ & $0{,}10$ & $0{,}12$\\\hline
\end{tabular}\end{center}
}{
\smallskip\noindent\textbf{Решение}\par
Из закона Гука 

\[k= \frac{F_{\text{упр}}}{x} = \frac{4}{0{,}04} =100\,\text{Н}/\text{м}.\]

Ответ: $100$.
}

% Source sdamgia-46762: https://phys-ege.sdamgia.ru/problem?id=46762
\BankTask{sdamgia-46762}{
\begin{center}\begin{tikzpicture}[book diagram,>=stealth]
\draw(0,0)--(5,0);\draw(0,0)--(0,1.6)(5,0)--(5,1.6);
\foreach \x in{0,.2,...,5}\draw(\x,0)--++(-.12,-.12);
\draw[fill=white](2,.05)rectangle(3,.95);\node at(2.5,.5){$M$};
\draw(0,.55)--(.2,.55)--(.35,.85)--(.55,.25)--(.75,.85)--(.95,.25)--(1.15,.85)--(1.35,.25)--(1.55,.85)--(1.75,.25)--(1.9,.55)--(2,.55);
\draw(3,.55)--(3.1,.55)--(3.25,.85)--(3.45,.25)--(3.65,.85)--(3.85,.25)--(4.05,.85)--(4.25,.25)--(4.45,.85)--(4.65,.25)--(4.8,.55)--(5,.55);
\node[above]at(1,1){$k_1$};\node[above]at(4,1){$k_2$};
\end{tikzpicture}\end{center}Кубик массой $100\,\text{г}$ покоится на гладком горизонтальном столе, сжатый с боков пружинами (см. рис.). Оси пружин совпадают. Жёсткость правой пружины $k_{2}$ = $600\,\text{Н}$/м. Левая пружина действует на кубик с силой $12\,\text{Н}$. На сколько сжата правая пружина? (В сантиметрах.)
}{
\smallskip\noindent\textbf{Решение}\par
Так как кубик покоится, то силы упругости, действующие на него со стороны каждой из пружин равны по модулю. Тогда из закона Гука найдем деформацию правой пружины:

\[x_2 = \frac{F_2}{k_2} = \frac{12}{600} = 0{,}02\,\text{м} = 2\,\text{см}.\]

Ответ: $2$.
}

% Source sdamgia-35249: https://phys-ege.sdamgia.ru/problem?id=35249
\BankTask{sdamgia-35249}{
Груз массой $200\,\text{г}$ подвешен на пружине жесткостью $100\,\text{Н}$/м к потолку лифта. Лифт равноускорено движется вниз, набирая скорость. Каково ускорение лифта, если удлинение пружины постоянно и равно $1{,}5\,\text{см}$?
}{
\smallskip\noindent\textbf{Решение}\par
Когда лифт покоился, на груз действовали сила тяжести mg и сила упругости пружины, которая по закону Гука равна $F_{\text{упр}}=kx.$ Поскольку груз находился в покое, то выполнялось равенство $mg=kx.$

При движении лифта вниз груз покоится относительно лифта, поэтому по второму закону Ньютона 

\[ma=mg - F_{\text{упр}};\]

\[ma=mg - kx.\]

Тогда ускорение груза равно:

\[a= g - \frac{kx}{m} = 10 - \frac{100 \cdot  0{,}015}{0{,}2} =2{,}5\,\text{м/с} ^2 .\]

Ответ: $2{,}5\,\text{м/с} ^2 .$
}

% Source sdamgia-3720: https://phys-ege.sdamgia.ru/problem?id=3720
\BankTask{sdamgia-3720}{
\begin{center}\begin{tikzpicture}[task graph,x=.75cm,y=.32cm]
\draw[task grid,xstep=.5,ystep=1] (0,0) grid (4.5,9);\draw[task axis] (0,0)--(4.9,0) node[right] {$x$, см};\draw[task axis] (0,0)--(0,9.6) node[above] {$F$, Н};\draw[task curve] (0,0)--(4.5,9);
\foreach \x in {1,2,3,4}{\node[below] at (\x,0) {$\x$};}\foreach \y in {2,4,6}{\node[left] at (0,\y) {$\y$};}\node[below left] at (0,0) {$0$};
\end{tikzpicture}\end{center}На графике приведена зависимость модуля силы упругости $F$ растянутой пружины от величины ее растяжения $x.$ Найти период свободных колебаний груза массой $0{,}5\,\text{кг}$, подвешенного на этой пружине.
}{
\smallskip\noindent\textbf{Решение}\par
Определим жесткость пружины. Согласно закону Гука, сила упругости пропорциональна удлинению пружины $F=kx,$ что и подтверждает данный график. Используя любую точку на графике для жесткости пружины, имеем: 

$k= \frac{F}{x} = \frac{4\,\text{Н}}{0{,}02\,\text{м}} =200\,\text{Н/м}.$

Период свободных колебаний пружинного маятника связан с жесткостью пружины и массой груза соотношением:

$T=2 \pi \sqrt{\frac{m}{k}} =2 \cdot 3{,}14 \cdot \sqrt{\frac{0{,}5\,\text{кг}}{200\,\text{Н/м}}} =0{,}314\,\text{с}\approx0{,}3\,\text{с}.$

Ответ: $0{,}3\,\text{с}$.
}

% Source sdamgia-50848: https://phys-ege.sdamgia.ru/problem?id=50848
\BankTask{sdamgia-50848}{
\begin{center}\begin{tikzpicture}[task graph]
\draw[task curve] (0,0)--(5,0);\draw[task curve] (0,0)--(0,2);
\foreach \x in {.1,.35,...,4.9}{\draw[task tick] (\x,0)--++(-.18,-.18);}\foreach \y in {.15,.4,...,1.9}{\draw[task tick] (0,\y)--++(-.18,-.18);}
\draw[task curve] (0,1)--(.5,1)--(.65,1.65)--(.95,.4)--(1.25,1.65)--(1.55,.4)--(1.85,1.65)--(2.15,.4)--(2.45,1.65)--(2.75,.4)--(2.9,1)--(3.1,1);
\draw[task curve,fill=orange!25] (3.1,.55) rectangle (4.8,1.7);\draw[task curve,fill=white] (3.45,.25) circle (.25);\draw[task curve,fill=white] (4.45,.25) circle (.25);
\end{tikzpicture}\end{center}Тележка массой $0{,}5\,\text{кг}$, прикрепленная к горизонтальной пружине жёсткостью $200\,\text{Н/м}$, совершает свободные гармонические колебания (см. рис.). Максимальная скорость тележки равна $3\,\text{м}/\text{с}$. Какова амплитуда колебаний тележки? Массой колёс можно пренебречь.
}{
\smallskip\noindent\textbf{Решение}\par
При гармонических колебаниях выполняется закон сохранения энергии: $\frac{kA ^2}{2} = \frac{m v _{\max} ^2}{2} ,$ откуда найдем амплитуду колебаний тележки:

$A = v _{\max} \sqrt{\frac{m}{k}} = 3 \cdot \sqrt{\frac{0{,}5}{200}} = 0{,}15\,\text{м} = 15\,\text{см}.$

Ответ: $15\,\text{см}$.
}


% Source 909F5c: https://ege.fipi.ru/bank/questions.php?search=1&proj=BA1F39653304A5B041B656915DC36B38&qid=909F5c
\BankTask{909F5c}{
На рисунке представлен график зависимости модуля силы упругости $F$ пружины от её удлинения $x$. Чему равна жёсткость пружины (в $\text{Н}/\text{м}$)?
\begin{center}
\begin{tikzpicture}[x=0.68cm,y=0.68cm,>=stealth]
\draw[step=1,black,thin] (0,0) grid (4,4);
\draw[->] (0,0)--(4.65,0) node[right] {$x,\,10^{-2}\,\text{м}$};
\draw[->] (0,0)--(0,4.75) node[left] {$F,\,\text{Н}$};
\foreach \n in {1,2,3,4}{\node[below] at (\n,0) {$\n$};\node[left] at (0,\n) {$\n$};}
\node[below left] at (0,0) {$0$};
\draw[thick] (0,0)--(4,4);
\end{tikzpicture}
\end{center}
}{
\smallskip\noindent\textbf{Решение от GPT}\par
По графику при $x=4\cdot10^{-2}\,\text{м}$ сила упругости равна $F=4\,\text{Н}$. По закону Гука\[k=\frac{F}{x}=\frac{4}{4\cdot10^{-2}}=100\,\text{Н}/\text{м}.\] Ответ: $100$.
}

% Source 26769F: https://ege.fipi.ru/bank/questions.php?search=1&proj=BA1F39653304A5B041B656915DC36B38&qid=26769F
\BankTask{26769F}{
Под действием силы величиной $6\,\text{Н}$ пружина удлинилась на $3\,\text{см}$. Чему равна величина силы, под действием которой удлинение этой пружины составит $4{,}5\,\text{см}$ (в $\text{Н}$)?
}{
\smallskip\noindent\textbf{Решение от GPT}\par
Для одной пружины по закону Гука $F_1=kx_1$ и $F_2=kx_2$. Следовательно,\[F_2=F_1\frac{x_2}{x_1}=6\cdot\frac{4{,}5}{3}=9\,\text{Н}.\] Ответ: $9$.
}

% Source 89A46D: https://ege.fipi.ru/bank/questions.php?search=1&proj=BA1F39653304A5B041B656915DC36B38&qid=89A46D
\BankTask{89A46D}{
Груз массой $200\,\text{г}$ подвешен на пружине к потолку неподвижного лифта. Лифт начинает двигаться и в течение $2\,\text{с}$ равноускоренно опускается вниз на расстояние $5\,\text{м}$. Каково удлинение пружины при опускании лифта, если её жёсткость $100\,\text{Н}/\text{м}$? Движение груза также считать равноускоренным, возникновением колебаний пренебречь. Удлинение выразите в $\text{см}$.
}{
\smallskip\noindent\textbf{Решение от GPT}\par
Лифт начинает движение из состояния покоя, поэтому $a=2s/t^2=2\cdot5/2^2=2{,}5\,\text{м}/\text{с}^{2}$. Ускорение направлено вниз. При $g=10\,\text{м}/\text{с}^{2}$ по второму закону Ньютона $mg-kx=ma$. Тогда\[x=\frac{m(g-a)}{k}=\frac{0{,}2(10-2{,}5)}{100}=0{,}015\,\text{м}=1{,}5\,\text{см}.\] Ответ: $1{,}5$.
}

% Source C97E5E: https://ege.fipi.ru/bank/questions.php?search=1&proj=BA1F39653304A5B041B656915DC36B38&qid=C97E5E
\BankTask{C97E5E}{
При исследовании зависимости модуля силы упругости $F_{\text{упр}}$ от удлинения пружины были получены следующие данные. Результат выразите в $\text{Н}/\text{м}$.
\begin{center}\begin{tabular}{|c|c|c|c|c|}\hline
$F_{\text{упр}},\,\text{Н}$ & $2{,}5$ & $5{,}0$ & $10{,}0$ & $12{,}5$ \\\hline
$\Delta x,\,\text{м}$ & $0{,}01$ & $0{,}02$ & $0{,}04$ & $0{,}05$ \\\hline
\end{tabular}\end{center}
Определите по результатам исследования жёсткость пружины.
}{
\smallskip\noindent\textbf{Решение от GPT}\par
По закону Гука $k=F/\Delta x=2{,}5/0{,}01=250\,\text{Н}/\text{м}$. Остальные столбцы дают то же отношение.
}

% Source 1ccc46: https://ege.fipi.ru/bank/questions.php?search=1&proj=BA1F39653304A5B041B656915DC36B38&qid=1ccc46
\BankTask{1ccc46}{
Под действием силы величиной $9\,\text{Н}$ пружина удлинилась на $4{,}5\,\text{см}$. Каким станет удлинение пружины под действием силы, равной по модулю $6\,\text{Н}$? Результат выразите в $\text{см}$.
}{
\smallskip\noindent\textbf{Решение от GPT}\par
По закону Гука $x_2/x_1=F_2/F_1$. Тогда $x_2=4{,}5\cdot6/9=3\,\text{см}$.
}

% Source BE8E13: https://ege.fipi.ru/bank/questions.php?search=1&proj=BA1F39653304A5B041B656915DC36B38&qid=BE8E13
\BankTask{BE8E13}{
Один конец лёгкой пружины жёсткостью $k$ прикреплён к бруску, а другой закреплён неподвижно. Брусок скользит по горизонтальной направляющей так, что его координата изменяется со временем по закону $x(t)=A\sin\omega t$.
Установите соответствие между физическими величинами и формулами, выражающими их изменения во времени.\TaskMatchingColumns[.53\linewidth]{Физические величины}{
А) проекция $F_x(t)$ равнодействующей силы на ось $x$\par
Б) потенциальная энергия пружины $E_{\text{П}}(t)$
}{Формулы}{
1) $kA^2\sin\omega t$\par
2) $-kA\sin\omega t$\par
3) $-kA^2\sin^2\omega t$\par
4) $\frac{kA^2}{2}\sin^2\omega t$\par
}\begin{center}\begin{tabular}{|c|c|}\hline А & Б \\\hline\rule{0pt}{1.2em}\hspace{1cm} & \hspace{1cm} \\\hline\end{tabular}\end{center}
}{
\smallskip\noindent\textbf{Решение от GPT}\par
По закону Гука $F_x=-kx=-kA\sin\omega t$ (формула $2$). Энергия $E_{\text{П}}=kx^2/2=(kA^2/2)\sin^2\omega t$ (формула $4$). Ответ: $24$.
}

% Source 46D7FD: https://ege.fipi.ru/bank/questions.php?search=1&proj=BA1F39653304A5B041B656915DC36B38&qid=46D7FD
\BankTask{46D7FD}{
Необходимо собрать экспериментальную установку, с помощью которой можно определить жёсткость резинового жгута. Для этого школьник взял штатив с лапкой и резиновый жгут. Какие два предмета из приведённого ниже перечня оборудования необходимо дополнительно использовать для проведения этого эксперимента?
\begin{enumerate}
\item[1)] мензурка
\item[2)] линейка
\item[3)] секундомер
\item[4)] стакан с водой
\item[5)] динамометр
\end{enumerate}
}{
\smallskip\noindent\textbf{Решение от GPT}\par
Линейкой (2) измеряют удлинение жгута \(\Delta l=l-l_0\), а динамометром (5) — растягивающую силу \(F\). В области выполнения закона Гука жёсткость равна \(k=F/\Delta l\). Мензурка (1), секундомер (3) и стакан с водой (4) для этих измерений не нужны. Ответ: 2, 5.
}

% Source 25774: https://phys-ege.sdamgia.ru/problem?id=25774
\BankTask{25774}{
Имеется недеформированная пружина длиной $L$  =  $30\,\text{см}$ и жесткостью $k$  =  $30\,\text{Н/м}$, груз массой $m$  =  $1\,\text{кг}$, а также вращающийся с частотой $\nu = 0{,}5\,\text{Гц}$ массивный диск. На каком минимальном расстоянии от центра диска можно положить на него груз, прикрепив его пружиной к центру диска, чтобы груз оставался неподвижным относительно диска? Коэффициент трения между грузом и диском $\mu=0{,}5.$ Размерами груза пренебречь. Сделайте схематический рисунок с указанием сил, действующих на груз.
Какие законы Вы использовали для описания движения бруска? Обоснуйте их применимость к данному случаю.
}{
\smallskip\noindent\textbf{Решение}\par
Обоснование. 
$1$.  Рассмотрим задачу в инерциальной системе отсчета, связанной с Землей. Будем считать эту систему отсчета инерциальной (ИСО). Тело описываем моделью материальной точки, так как его размеры малы по сравнению с диском. 
$2$.  При вращении диска на тело действуют потенциальная сила тяжести, сила реакции опоры, перпендикулярная опоре, потенциальная сила упругости, возникающая в упруго деформированной пружине и сила трения покоя. Работа силы тяжести и силы реакции опоры равны нулю. 
$3$.  Поскольку тело описывается моделью материальной точки, в ИСО применим второй закон Ньютона. Равнодействующая сила направлена к центру, следовательно, тело движется с центростремительным ускорением. Для материальной точки применимы законы равномерного движения по окружности.
$4$.  При малых деформациях их можно считать упругими, поэтому для пружины верен закон Гука.
$5$.  Брусок покоится относительно диска, имеющего шероховатую поверхность. Поэтому на него действует максимальная сила трения покоя, равная силе трения скольжения $F_{\,\text{тр}}=\mu N.$

Перейдем к решению. \begin{center}\begin{tikzpicture}[task graph,>=stealth]\draw(.1,0)--(5,0)--(5,-.4)--(.1,-.4);\draw[dashed](.1,0)--(-.85,0)--(-.85,-.4)--(.1,-.4);\draw(-.4,-.4)--(.1,-.4)--(.1,1.8);\draw[very thick](.1,-.4)--(.1,1.8);\draw[fill=gray!15](3,.05)rectangle(4.2,.85);\draw(.1,.45)--(.3,.45);\draw plot[domain=0:1,samples=100]({.3+2.7*\x},{.45+.15*sin(2880*\x)});\node[below right]at(5.2,-1.1){$O$};\draw[->](5.2,-1.1)--(5.2,2.1)node[above]{$y$};\draw[->](5.2,-1.1)--(1.2,-1.1)node[left]{$x$};\draw[->](.1,1.8)..controls(-.5,2.3)and(.6,2.3)..(.1,1.8)node[above left]{$\nu$};\draw[->,thick](3.6,.45)--(3.6,2.3)node[above]{$\vec N$};\draw[->,thick](3.6,.45)--(3.6,-.9)node[right]{$m\vec g$};\draw[->,thick](3.6,.45)--(5,.45)node[right]{$\vec F_{\text{упр}}$};\draw[->,thick](3.6,.45)--(1.9,.45)node[above left]{$\vec F_{\text{тр}}$};\draw[->,thick](2.6,2.3)--(1.4,2.3)node[above left]{$\vec a_{\text{ц}}$};\end{tikzpicture}\end{center} 
$1$.  На груз действуют четыре силы: сила тяжести $m\vec g,$ нормальная составляющая силы реакции опоры $\overrightarrow{N},$ сила упругости $\overrightarrow{F}_{\,\text{упр}}$ и сила трения $\overrightarrow{F}_{\,\text{тр}}.$
Выясним, в каком состоянии будет находиться пружина. Если грузик находится на расстоянии $L$ от центра, то равнодействующая сила равна $ma=4 \pi ^2 \nu ^2 Lm=3$$H$. Максимальная сила трения равна $F_{\,\text{т}} =\mu mg =5$$H$. Значит, пружина должна быть сжата, сила упругости направлена вправо (см. рис.), а сила трения  — влево. Поскольку пружина максимально сжата, а груз при этом находится в покое относительно диска, сила трения покоя принимает максимальное значение: $F_{\,\text{тр}}=\mu N.$
$2$.  Второй закон Ньютона в проекциях на оси инерциальной системы отсчета XOY имеет вид: $OX: ma_{\,\text{ц}}=F_{\,\text{тр}} - F_{\,\text{упр}};$ $OY: 0=N - mg.$ Сила упругости определяется законом Гука: $F_{\,\text{упр}}=k\Delta x.$ Центростремительное ускорение груза выражается формулой $a_{\,\text{ц}}= \frac{V ^2}{R_{\min}} =4 \pi ^2 \nu ^2 R_{\min}.$
$3$.  Из формул, приведенных выше, получаем: $4 \pi ^2 \nu ^2 R_{\min}m=\mu mg - k\Delta x,$ где $R_{\min}=L - \Delta x.$ Окончательно получим:

\[R_{\min}= \frac{kL - \mu mg}{k - 4 \pi ^2 \nu ^2 m} = \frac{30 \cdot 0{,}3 - 0{,}5 \cdot 1 \cdot 10}{30 - 4 \cdot 3{,}14 ^2 \cdot 0{,}5 ^2 \cdot 1} \approx 0{,}2 \,\text{м}.\]

Ответ: $R_{\min}\approx 0{,}2 \,\text{м}.$
}

% Source 25407: https://phys-ege.sdamgia.ru/problem?id=25407
\BankTask{25407}{
\begin{center}\begin{tikzpicture}[task graph,>=stealth]
\draw(0,3.5)--(4,3.5);\foreach \x in {0,.2,...,4}{\draw(\x,3.5)--++(.12,.12);}
\draw[decorate,decoration={zigzag,segment length=4mm,amplitude=1.4mm}](.4,3.4)--(.4,1.6);
\draw[decorate,decoration={zigzag,segment length=4mm,amplitude=1.4mm}](3.6,3.4)--(3.6,1.6);
\draw(.4,3.5)--(.4,3.4)(3.6,3.5)--(3.6,3.4)(.4,1.6)--(.4,1.5)(3.6,1.6)--(3.6,1.5);
\draw(.2,1.3)rectangle(3.8,1.5);\draw(3,1.3)--(3,.8);\draw(2.7,.2)rectangle(3.3,.8);
\end{tikzpicture}\end{center}К двум вертикально расположенным пружинам одинаковой длины подвесили однородный стержень длиной $L=30\,\text{см}.$ Если к этому стержню подвесить груз массой $m=3\,\text{кг}$ на расстоянии $d=5\,\text{см}$ от правой пружины, то стержень будет расположен горизонтально, и растяжения обеих пружин будут одинаковы (см. рис.). Жесткость левой пружины в $2$ раза меньше, чем правой. Чему равна масса стержня $M$? Сделайте рисунок с указанием используемых в решении сил.
Какие законы Вы использовали для описания равновесия системы? Обоснуйте их применимость к данному случаю.
}{
\smallskip\noindent\textbf{Решение}\par
Обоснование
$1$.  Рассмотрим задачу в системе отсчета, связанной с Землей. Будем считать эту систему отсчета инерциальной (ИСО).
$2$.  Описываем стержень моделью твердого тела (формы и размеры тела неизменны, расстояние между любыми двумя точками тела остается неизменным).
$3$.  Любое движение твердого тела является суперпозицией поступательного и вращательного движений. Поэтому условий равновесия твердого тела в ИСО ровно два: одно для поступательного движения, другое - для вращательного движения.
$4$.  Сумма приложенных к твердому телу внешних сил равна нулю (условие равновесия твердого тела относительно поступательного движения). Поэтому сумма моментов этих сил относительно любых двух параллельных осей одна и та же. Для удобства выберем ось, проходящую перпендикулярно плоскости рисунка через центр масс стержня (точку О).
Решение
\begin{center}\begin{tikzpicture}[task graph,>=stealth]\draw(-.3,3)--(5.3,3);\draw(0,3)--(0,2.8);\draw plot[domain=0:1,samples=80]({0+.1*sin(2520*\x)},{2.8-1.5*\x});\draw(4.8,3)--(4.8,2.8);\draw plot[domain=0:1,samples=80]({4.8+.1*sin(2520*\x)},{2.8-1.5*\x});\draw[fill=gray!15](-.15,1.1)rectangle(4.95,1.3);\draw(3.7,1.1)--(3.7,.65);\draw[fill=gray!15](3.3,0)rectangle(4.1,.65);\node[above]at(2.4,1.4){$A$};\node[right]at(4.95,1.25){$B$};\draw[->,thick](0,1.3)--(0,2.3)node[above left]{$\vec F_{\text{у}1}$};\draw[->,thick](4.8,1.3)--(4.8,2.7)node[right]{$\vec F_{\text{у}2}$};\draw[->,thick](2.4,1.1)--(2.4,-.3)node[below left]{$M\vec g$};\draw[->,thick](3.7,.3)--(3.7,-.7)node[right]{$m\vec g$};\end{tikzpicture}\end{center}
$1$.  Укажем на рисунке силы действующие на стержень. Приравняем моменты сил, действующих на стержень, относительно центра стержня, то есть точки А:

\[F_{\,\text{у}1} \cdot \frac{L}{2} + mg \cdot ( \frac{L}{2} - d ) =F_{\,\text{у}2} \cdot \frac{L}{2} .\]

Учтем, что стержень расположен горизонтально, то есть удлинения пружин равны, а также, что жесткость правой пружины в два раза больше левой:

\[kx \cdot \frac{L}{2} + mg \cdot ( \frac{L}{2} - d ) =2kx \cdot \frac{L}{2} \Rightarrow F_{\,\text{у}1}=kx= \frac{mg ( L - 2d )}{L} .\]

$2$.  Приравняем моменты сил, действующих на стержень, относительно точки Б, которая находится в месте крепления правой пружины:

\[kx \cdot L=Mg \cdot \frac{L}{2} + mg \cdot d.\]

$3$.  Найдем массу стержня:

\[M= \frac{kx \cdot L - mg \cdot d}{g \cdot \frac{L}{2}} = \frac{mg ( L - 2d ) - mg \cdot d}{g \cdot \frac{L}{2}} = \frac{2m ( L - 3d )}{L} = \frac{2 \cdot 3 \cdot ( 30 - 15 )}{30} =3\,\text{кг}.\]

Ответ: $M= \frac{2m ( L - 3d )}{L} =3\,\text{кг}.$
}

\end{document}
